%% ── 11. LIMITACIONES Y RIESGOS ───────────────────────────────────────────────
\chapter{Limitaciones y Riesgos}

Este capítulo documenta las limitaciones identificadas durante el desarrollo
y evaluación del MVP, organizadas por categoría, y presenta la matriz de
riesgos residuales con probabilidad, impacto y estrategia de mitigación.
Reconocer estas limitaciones no invalida los resultados del sistema, sino
que define con precisión el marco dentro del cual deben interpretarse
y orienta las prioridades del roadmap del Capítulo 5.

\section*{Limitaciones Técnicas}

\begin{enumerate}

    \item \textbf{Variabilidad de formato en PDFs:}
    Los contratos de concesión peruanos presentan alta heterogeneidad
    estructural: numeraciones mixtas (arábiga + romana + literal),
    tablas de múltiples columnas dentro de cláusulas, y anexos técnicos
    en formato de planilla. El módulo de segmentación fue diseñado para el
    patrón estructural del Contrato A. Contratos con:
    \begin{itemize}
        \item OCR defectuoso o texto escaneado sin normalización,
        \item tablas de datos cuantitativos en lugar de texto narrativo,
        \item numeración no jerárquica (por ejemplo, cláusulas identificadas
              solo por letra sin número de sección),
    \end{itemize}
    pueden requerir ajustes manuales en las heurísticas de segmentación
    o reducir la precisión de la detección de límites entre cláusulas.
    \textit{Mitigación parcial:} el módulo anti-TOC y el filtro de longitud
    mínima ($>200$ caracteres) reducen el impacto, pero no lo eliminan.

    \item \textbf{Límite de ventana de contexto:}
    Los LLMs tienen una ventana de tokens finita. Cláusulas extremadamente
    extensas ($> 8{,}000$ tokens) se truncan antes de enviarse al agente,
    con riesgo de perder información relevante al final del fragmento.
    \textit{Mitigación parcial:} truncamiento adaptativo al 80\% del límite
    preservando la cabecera de la cláusula; el Reranker prioriza el fragmento
    más relevante. El problema persiste en secciones con listas de obligaciones
    muy largas (algunos anexos del VIC PTAR superan los 6{,}000 tokens
    por cláusula).

    \item \textbf{Variabilidad inherente a los LLMs generativos:}
    Con temperatura 0.1, las ejecuciones sobre el mismo contrato producen
    resultados muy similares pero no idénticos. En pruebas de repetibilidad
    (3 ejecuciones consecutivas sobre el mismo input), se observaron
    variaciones de $\pm 3$--$5$ hallazgos, principalmente en hallazgos
    de severidad BAJA y en la redacción de explicaciones textuales.
    \textit{Mitigación parcial:} esquemas Pydantic con validación estricta
    de tipos; deduplicación por similitud coseno ($> 0.90$). La
    variabilidad en el conteo total de hallazgos no supera el 4\%.

    \item \textbf{Concurrencia limitada a un análisis simultáneo:}
    El MVP usa \texttt{asyncio.Semaphore(1)} para garantizar que solo
    un análisis completo se ejecute a la vez. Esto evita la saturación
    de la cuota de Vertex AI y la sobrecarga de memoria en Cloud Run,
    pero impide que dos usuarios realicen análisis de forma estrictamente
    simultánea. El sistema de cola por usuario (máx.\ 5 trabajos)
    mitiga parcialmente el impacto, pero en escenarios de uso intensivo
    (varios auditores simultáneos) la espera puede superar 75 minutos.

    \item \textbf{Dependencia de conectividad con APIs externas:}
    El pipeline requiere conexión activa a Vertex AI (LLM + embeddings)
    y a Cohere (Reranker). Un corte o degradación del servicio de
    cualquiera de estos proveedores detiene el análisis. No existe modo
    de operación offline ni fallback local de inferencia en el MVP.

\end{enumerate}

\section*{Limitaciones de Datos}

\begin{enumerate}

    \item \textbf{Cobertura normativa parcial:}
    La base normativa vectorizada incluye: D.Leg.\ 1362 y su Reglamento,
    guías institucionales de ProInversión, y normativa SUNASS sectorial.
    No cubre:
    \begin{itemize}
        \item Normativa de otros sectores (energía, telecomunicaciones, vial).
        \item Resoluciones administrativas recientes (post-2024).
        \item Precedentes arbitrales relevantes.
        \item Contratos modelo o bases estándar de licitación.
    \end{itemize}
    La ausencia de estas fuentes puede provocar falsos negativos en el
    módulo \texttt{OMISION\_NORMATIVA}: el sistema no detecta una omisión
    si la norma infringida no está vectorizada en la base RAG.

    \item \textbf{Dominio restringido a contratos APP peruanos:}
    El sistema fue diseñado, entrenado (en términos de prompts y base
    normativa) y validado exclusivamente sobre contratos de concesión APP
    bajo el marco legal peruano. La aplicación a:
    \begin{itemize}
        \item contratos de otro tipo (EPC, O\&M, suministro, consultoría),
        \item contratos bajo legislación extranjera o marcos bilingües,
        \item documentos no contractuales (bases de licitación, adendas),
    \end{itemize}
    requiere una revisión completa de prompts, base normativa y esquemas
    de clasificación de hallazgos.

    \item \textbf{Validación sobre un único contrato piloto:}
    Los indicadores de precisión ($\approx 87\%$) y recall ($\approx 94\%$)
    se calcularon sobre una muestra del Contrato A. No existe
    evidencia estadística de que estos valores se mantengan en contratos
    de otros sectores, años o niveles de complejidad. La generalización
    requiere validación cruzada sobre al menos 5--10 contratos distintos.

\end{enumerate}

\section*{Limitaciones Operativas}

\begin{enumerate}

    \item \textbf{Interfaz en etapa MVP:}
    El frontend Next.js es funcional pero no está diseñado para operación
    institucional de larga duración. Carece de:
    \begin{itemize}
        \item Sistema de roles granular (legal, técnico, ambiental, auditor,
              director) con permisos diferenciados.
        \item Auditoría de acciones de usuario con log inmutable.
        \item Soporte para contratos en idioma inglés o formatos bilingües.
        \item Panel de administración completo para gestión de la base
              normativa.
    \end{itemize}

    \item \textbf{Sin integración con sistemas institucionales:}
    El MVP opera de forma autónoma. No está conectado a repositorios
    documentales de ProInversión, sistemas de gestión contractual, flujos
    de aprobación electrónica ni plataformas de firma digital. La carga
    de contratos es manual (upload por usuario) y los reportes se
    descargan individualmente.

    \item \textbf{Costo operativo por análisis:}
    El costo estimado por análisis completo con Gemini~2.5~Pro es de
    $\$50$--$\$100$, compuesto principalmente por tokens LLM (Vertex AI)
    y embeddings. A escala institucional (30--50 contratos/año), el costo
    anual estimado sería de $\$1{,}500$--$\$5{,}000$, que aunque es muy
    inferior al costo manual, puede representar una barrera en entidades
    con presupuesto de TI limitado.

\end{enumerate}

\section*{Limitaciones de Negocio}

\begin{enumerate}

    \item \textbf{El sistema es apoyo, no sustituto del especialista:}
    El $\approx 13\%$ de falsos positivos y los $\approx 8$ falsos
    negativos estimados confirman que el MVP actúa correctamente como
    herramienta de \textit{screening} previo. Los hallazgos de severidad
    ALTA requieren validación experta antes de constituir observaciones
    formales, recomendaciones o actos administrativos. Usar el reporte
    del sistema como insumo directo para decisiones legales sin revisión
    experta constituye un uso incorrecto del producto.

    \item \textbf{Cuota de modelos Claude pendiente:}
    La solicitud de cuota en Vertex AI para Claude Sonnet~4.6 y
    Claude Opus~4.6 fue rechazada durante el periodo de evaluación.
    La evaluación comparativa con modelos Anthropic está pendiente y
    podría modificar las conclusiones sobre el modelo óptimo de producción.

    \item \textbf{Resistencia al cambio institucional:}
    La adopción de herramientas de IA en entidades públicas enfrenta
    barreras culturales y regulatorias: desconfianza en los resultados,
    ausencia de marcos legales que regulen el uso de IA en auditoría
    contractual, y procesos de aprobación tecnológica lentos. Sin un
    programa de gestión del cambio, el producto técnico puede no
    traducirse en adopción real.

\end{enumerate}

\section*{Matriz de Riesgos Residuales}

Los riesgos residuales son aquellos que persisten después de las medidas
de mitigación implementadas en el MVP. Se clasifican por probabilidad
e impacto en una escala Alta/Media/Baja:

\begin{longtable}{p{3.8cm} p{1.5cm} p{1.5cm} p{6.5cm}}
\caption{Matriz de riesgos residuales del MVP}
\label{tab:riesgos} \\
\toprule
\textbf{Riesgo} & \textbf{Prob.} & \textbf{Impacto} & \textbf{Mitigación} \\
\midrule
\endfirsthead
\multicolumn{4}{l}{\small\itshape Continuación de la Tabla~\ref{tab:riesgos}} \\[4pt]
\toprule
\textbf{Riesgo} & \textbf{Prob.} & \textbf{Impacto} & \textbf{Mitigación} \\
\midrule
\endhead
\midrule
\multicolumn{4}{r}{\small\itshape Continúa en la siguiente página} \\
\endfoot
\bottomrule
\endlastfoot
Deprecación de Gemini~2.5~Pro por Google &
    Baja & Alto &
    Capa de abstracción \texttt{provider.py}; variable \texttt{LLM\_MODEL}
    configurable; fallback documentado. \\[8pt]
Cambio de API Vertex AI o Cohere &
    Baja & Medio &
    Versionado de clientes en \texttt{requirements.txt};
    pruebas de integración ante cada actualización. \\[8pt]
Costo de inferencia excesivo a escala &
    Media & Medio &
    FinOps IA: caché de embeddings normativos; reducción de contexto;
    evaluación de modelos más eficientes. \\[8pt]
Fallo de Cloud SQL o Cloud Run &
    Baja & Alto &
    Cloud SQL con backups automáticos diarios; Cloud Run con health
    checks; rollback automático en CI/CD. \\[8pt]
Resistencia institucional a adopción &
    Media & Alto &
    Workshops de adopción; reporte ejecutivo legible por no técnicos;
    demostración en contrato real con hallazgos verificados. \\[8pt]
Contrato con estructura no convencional &
    Media & Medio &
    Validación previa de segmentación; modo de revisión manual del
    índice de secciones antes del análisis completo. \\[8pt]
Violación de confidencialidad documental &
    Baja & Alto &
    Contratos almacenados con cifrado en Cloud Storage; acceso por
    signed URLs de duración limitada; políticas IAM por servicio. \\[8pt]
Pérdida de acceso a Vertex AI (cuota agotada) &
    Media & Alto &
    Monitoreo de uso en Vertex AI; alertas automáticas al 80\% de cuota;
    sistema de cola evita ráfagas de solicitudes. \\
\end{longtable}

\section*{Mapa de Mitigación en el Roadmap}

Todas las limitaciones identificadas tienen una contrapartida en las fases
del roadmap descrito en el Capítulo 5:

\begin{table}[H]
\centering
\caption{Limitaciones y su tratamiento en el roadmap de siguiente fase}
\label{tab:limitaciones_roadmap}
\begin{tabular}{p{5cm} p{2.5cm} p{5cm}}
\toprule
\textbf{Limitación} & \textbf{Fase} & \textbf{Acción prevista} \\
\midrule
Variabilidad de formato PDF     & Fase 1 & Heurísticas adicionales por sector \\
Cobertura normativa parcial     & Fase 1 & Ampliar base; pipeline de actualización \\
Concurrencia limitada           & Fase 3 & Colas distribuidas; Docker/Kubernetes \\
Interfaz MVP sin roles          & Fase 2 & Autenticación por roles institucionales \\
Sin integración institucional   & Fase 2 & APIs documentales; conector repositorios \\
Costo operativo a escala        & Fase 3 & FinOps IA; caché; modelos eficientes \\
Validación sobre 1 contrato     & Fase 1 & Benchmarks en 5+ contratos de sectores distintos \\
Resistencia institucional       & Fase 4 & Capacitación; workshops; certificación \\
\bottomrule
\end{tabular}
\end{table}

Las limitaciones actuales del MVP son tratables y están contempladas en el
plan de desarrollo. Ninguna constituye un impedimento estructural para la
evolución del sistema hacia una plataforma institucional de auditoría
contractual.